%----------------------------------------------------------------------------------------
%	COVER LETTER
%----------------------------------------------------------------------------------------

% To remove the cover letter, comment out this entire block
\clearpage

 
\recipient{}{} % Letter recipient
\date{\today} % Letter date
\opening{\textbf{Summer Funding Reserch Plan}} % Opening greeting
\closing{} % Closing phrase
\enclosure[Attached]{Curriculum Vit\ae{}} % List of enclosed documents


% \makelettertitle
\textbf{Summer Funding Plan}\par
I intend to use the Hamilton Award to enable full-time commitment to research this summer. Doing so will allow me to focus the entirety of my energies to research, speeding progress by enabling me to advance my research goals more quickly. \par 

Tissues and individual cells are poroelastic, consisting of a background matrix suffused by a solvent carrying solutes of varying sizes. In vivo, materials undergo dynamic stresses, causing deformations which induce advection and size-dependent transport. Solute distribution in the nucleus, a heterogeneous poroelastic environment often subjected to dynamic strains, drives gene transcription. Experiments examining the macroscale dynamics of poroelastic media demonstrate enhanced concentration of large solutes in cyclically strained hydrogels, but there remains a dearth of studies probing the intricate interactions between solutes, pores, stress gradients, and microfluidics during dynamic strains, thus the roots of these emergent behaviors remain unclear. Diffusive transport of large solutes in poroelastic media is anomalous, strain-dependent, and temporally anisotropic, necessitating high spatiotemporal resolution. \par

To date, my work has provided the first nanoscopic, single-particle insight into the dynamics of these system. Both static and dynamic strains increase long-term localization and decrease diffusivity compared to unstrained controls. Strains also enhances the magnitude and duration of non-Gaussian diffusion. Dynamic strain induces asymmetric advective transport of large solutes during compression and relaxation, while also reducing long-term solute localization asymmetrically throughout their cycle by as much as threefold compared to static strains. \par

 % I am investigating the impact of strains on diffusive transport in poroelastic media. A spherical indenter mounted on a piezo-actuated atomic force microscope tip is used to exert a precise strain on a polyacrylamide gel substrate immersed in a solution of quantum dots in a watertight PDMS chamber. The Quantum dots are of similar size to the average pore size of the polyacrylamide gel, small enough to infiltrate the gel but large enough that they behave as large solutes. Highly inclined swept tile light sheet fluorescence microscopy is used to image individual quantum dots at hundreds of frames per second. These images are used to conduct single particle tracking; the tracking data is then used to analyze the diffusive motion of the quantum dots as a stochastic process, granting insight into how large solutes in biological poroelastic media are impacted by physiological strains \textit{in vivo}. The reductions in diffusivity found in other studies are due to greater confinement of the solutes, while a pumping effect observed in poroelastic media subjected to dynamic strains are caused by an asymmetry in transport induced by advective flows. This work provides the first single particle tracking of solutes in strained poroelastic media, a poorly understood topic with far-reaching implications. \par


Over the course of the summer, I will apply strains to cells to understand how strains impact the transport of transcription factors \textit{in situ}. To accomplish this goal, individual fluorescently labelled transcription factors must be imaged in a live cellular nuclei being subjected to well-characterized strains. 

To accomplish this, human cells will be cultured by collaborators with extensive experience. These cells will be treated with a non-toxic HALO-tag solution to fluorescently label transcription factors (TFs) within the cell nuclei. Our microscope operates with low phototoxicity inside a climate-controlled incubation system, allowing cell health to be maintained for extended experimental sessions. Precise strain will be applied during live imaging using a spherical indenter mounted on a piezo-actuated atomic force microscope cantilever. This strategy was developed and successfully applied in similar prior experiments involving a poroelastic gel, and should be adaptable to this experiment with minimal alteration.

High speed imaging fluorescently-labelled TFs will be conducted using light sheet fluorescence microscopy to acquire high sensitivity, high-framerate time series of the particles. The necessary microscopy techniques were developed and in the course of previous work; this system has imaged quantum dots at up to 800fps, and this task should be well within demonstrated capabilities. Particle tracking will be conducted using TrackMate, open-source state-of-the-art particle tracking software, and the data will be analyzed using MATLAB scripts. This experiment is likely to be very similar to a recent collaboration tracking HALO-tagged histone protein complexes in the cellular nuclei, and the lessons learned there can be applied here.\par

 As an extension, a unique in-house instrument which enables simultaneous atomic force microscopy and light sheet fluorescence microscopy could also be used to produce simultaneous particle tracking and rheological data to examine diffusive dynamics as a function of strain.\par

\newpage
